% !TEX program = xelatex
\documentclass[UTF8,11pt,a4paper]{ctexart}

\usepackage[a4paper,margin=2.2cm]{geometry}
\usepackage{amsmath,amssymb,amsfonts,mathtools,bm}
\usepackage{booktabs}
\usepackage{multirow}
\usepackage{array}
\usepackage{enumitem}
\usepackage{graphicx}
\usepackage{hyperref}
\usepackage{algorithm}
\usepackage{algpseudocode}
\usepackage{xcolor}

\hypersetup{colorlinks=true,linkcolor=blue,citecolor=blue,urlcolor=blue}
\setlist[itemize]{topsep=2pt,itemsep=2pt,parsep=0pt,leftmargin=1.6em}
\setlist[enumerate]{topsep=2pt,itemsep=2pt,parsep=0pt,leftmargin=1.8em}

\title{NOVA\_EDRL：面向非平稳联运调度的双层环境设计强化学习\\论文初稿（中文）}
\author{34959\_RL 项目组}
\date{2026-03-03}

\begin{document}
\maketitle

\begin{abstract}
针对非平稳联运调度中“训练分布与部署分布错位”导致的动作塌缩与泛化退化问题，本文提出一体化双层框架 NOVA\_EDRL。该框架在既有内层 RL+ALNS 求解器之外引入外层环境设计器，以“生成训练分布参数”作为显式决策变量，形成 Phase1 基线预训练、Phase2 冻结内层训练外层、Phase3 外内联合课程、Phase4 复用 Master implement 测试的四阶段闭环。方法上，本文的关键技术包括：1）外层动作空间严格参数化并移除 seed 动作维，避免伪相关优化；2）面向 Phase2 设计 hard-but-learnable 目标函数，在对抗强度、反塌缩增益与可学习性护栏之间实现可控平衡；3）保持 Phase3 原 EDRL 目标不变，保证与既有工程系统兼容；4）构建与 RARL、PLR/UED、CQL、CaDM 并列的多 SOTA 对照矩阵，并给出公式级映射与实现口径。本文给出完整数学定义、目标函数、动作空间、收敛切换规则及工程实现映射。最终结论留空，待正式实验填充。
\end{abstract}

\tableofcontents
\newpage

\section{引言与研究问题}
\subsection{背景}
联运调度系统在真实部署中面临多源非平稳扰动：拥堵强度、订单到达节奏、突发约束和结构性缺失同时变化。若仅在单一静态分布训练，策略在 OOD（out-of-distribution）条件下常出现：
\begin{itemize}
\item \textbf{动作塌缩}：二值动作长期退化为单一动作；
\item \textbf{策略脆弱}：测试回报对分布轻微漂移高度敏感；
\item \textbf{训练-测试错配}：训练样本覆盖不足，部署失败难以被训练目标直接纠正。
\end{itemize}

\subsection{问题形式化}
定义内层调度环境为马尔可夫决策过程
\begin{equation}
\mathcal{M}(E)=\langle \mathcal{S},\mathcal{A},P_E,r,\gamma \rangle,
\end{equation}
其中 \(E\) 为事件文件集合（由外层生成），\(\mathcal{A}=\{0,1\}\) 为二值动作空间。给定内层策略 \(\pi_\theta\)，其回报定义为
\begin{equation}
J(\theta;E)=\mathbb{E}_{\tau\sim (\pi_\theta,\mathcal{M}(E))}\!\left[\frac{1}{T}\sum_{t=1}^{T} r_t\right],\quad r_t\in\{0,1\}.
\end{equation}

外层策略 \(\pi_\psi\) 不直接控制运输动作，而是选择环境参数动作 \(a_g\)，并通过生成器 \(G\) 产生训练分布：
\begin{equation}
a_g \sim \pi_\psi(\cdot|s_t^g), \qquad E_t = G(a_g).
\end{equation}
目标是在可学习前提下最大化内层在多分布测试上的稳健性。

\subsection{本文贡献}
\begin{enumerate}
\item 提出 NOVA\_EDRL 四阶段一体化链路，并在 Master 中作为独立算法入口实现。
\item 给出 Phase2 新目标函数：外层在固定 \(\theta_{\text{phase1}}\) 下寻找“难但可学”的分布参数。
\item 保持 Phase3 目标函数与既有 EDRL 一致，确保历史对照与工程兼容性。
\item 建立多 SOTA 矩阵（RARL/PLR-UED/CQL/CaDM），形成可复现实验比较框架。
\end{enumerate}

\section{NOVA\_EDRL 总体框架}
\subsection{四阶段主流程}
当算法选择 \texttt{NOVA\_EDRL} 时，系统固定执行：
\begin{enumerate}
\item \textbf{Phase1（内层基线预训练）}：\texttt{stage-mode=train\_only}，训练得到 \(\theta_{\text{phase1}}\)；
\item \textbf{Phase2（外层训练，内层冻结）}：每轮内层都从同一 \(\theta_{\text{phase1}}\) 启动，外层更新；
\item \textbf{Phase3（联合训练）}：恢复内层跨轮继承更新（\(\theta_t\rightarrow\theta_{t+1}\)），外内课程协同；
\item \textbf{Phase4（实现复用测试）}：回到 Master implement，\texttt{stage-mode=eval\_only} 做统一评估。
\end{enumerate}

该流程对应实现：
\begin{itemize}
\item \texttt{codes/Dynamic\_master34959.py}
\item \texttt{codes/outer\_rl/run\_edrl\_pipeline.py}
\item \texttt{codes/outer\_rl/run\_edrl\_phase2.py}
\end{itemize}

\subsection{外层动作空间}
外层动作定义为
\begin{equation}
a_g=(\mu_A,\mu_B,p,n,\mathrm{pattern}).
\end{equation}
默认候选集合（可配置）：
\begin{align}
\mu_A,\mu_B &\in \mathcal{M}=\{10,30,60,90\},\ \mu_A\neq\mu_B, \\
p &\in \mathcal{P}=\{0.2,0.3,0.5,0.7,0.8\}, \\
n &\in \mathcal{N}=\{5,10,15,20,25,30\},\quad 3\le n\le 200, \\
\mathrm{pattern} &\in \mathcal{T}=\{\texttt{ab},\texttt{random\_mix}\}.
\end{align}
离散动作总数为
\begin{equation}
|\mathcal{A}_g| = |\mathcal{M}|(|\mathcal{M}|-1)|\mathcal{P}||\mathcal{N}||\mathcal{T}|.
\end{equation}

\subsection{seed 去动作化}
外层不再学习 seed。每轮 seed 由基础随机种子和迭代号确定性映射：
\begin{equation}
\mathrm{seed}_t = h(\mathrm{base\_seed},t),
\end{equation}
其中 \(h\) 为固定哈希映射。该策略减少无效搜索维度，避免“外层学习随机标签”的伪优化行为。

\subsection{外层策略实现族}
当前代码支持：
\begin{equation}
\texttt{policy\_mode}\in\{\texttt{random},\texttt{ucb},\texttt{ts},\texttt{pg},\texttt{rarl\_dqn}\}.
\end{equation}
关键更新规则如下。
\paragraph{UCB}
\begin{equation}
i_t=\arg\max_i\left(\bar Q_i + c\sqrt{\frac{\ln(1+N_t)}{1+n_i}}\right).
\end{equation}
\paragraph{Thompson Sampling}
\begin{equation}
\tilde Q_i \sim \mathcal{N}(\mu_i,\sigma_i^2),\qquad i_t=\arg\max_i \tilde Q_i.
\end{equation}
\paragraph{Policy-Gradient Bandit}
\begin{align}
\pi_\psi(i) &= \frac{\exp(z_i)}{\sum_j \exp(z_j)},\\
z_i &\leftarrow z_i + \alpha (R_t-b_t)\left(\mathbb{I}[i=i_t]-\pi_\psi(i)\right).
\end{align}
\paragraph{RARL-DQN 外层}
\begin{align}
Q_\psi(s_t^g,a_t^g) &\leftarrow Q_\psi(s_t^g,a_t^g) + \eta \delta_t,\\
\delta_t &= r_t^g + \gamma\max_a Q_{\bar\psi}(s_{t+1}^g,a)-Q_\psi(s_t^g,a_t^g).
\end{align}

\section{非稳态与 OOD 机理补充（结合既有研究稿）}
\subsection{从平稳 MDP 到上下文驱动非平稳过程}
标准平稳 MDP 写作：
\begin{equation}
\mathcal{M}=\langle \mathcal{S},\mathcal{A},P,R,\gamma\rangle.
\end{equation}
本系统存在事件阶段切换与分布漂移，转移与奖励更接近时变混合：
\begin{equation}
P_t(s'|s,a)=\sum_z P(s'|s,a,z)p_t(z|h_t),\qquad
R_t(s,a)=\sum_z R(s,a,z)p_t(z|h_t),
\end{equation}
其中 \(z_t\) 为潜在上下文（例如 phase、拥堵强度、事件生命周期），\(h_t\) 为历史。若 agent 仅看到压缩观测 \(o_t=\psi(s_t,z_t)\)，则问题本质上是 POMDP。

\subsection{Bellman 失配与 bootstrap 方差放大}
平稳情形下有
\begin{equation}
Q^\pi(s,a)=\mathbb{E}[r_t+\gamma \mathbb{E}_{a'\sim\pi}Q^\pi(s_{t+1},a')].
\end{equation}
非平稳下应写为时变版本
\begin{equation}
Q_t^\pi(s,a)=\mathbb{E}[r_t+\gamma \mathbb{E}_{a'\sim\pi}Q_{t+1}^\pi(s_{t+1},a')].
\end{equation}
单一参数网络拟合时变目标将导致目标噪声上升：
\begin{equation}
\mathrm{Var}[y_t]
=\mathrm{Var}[r_t]+\gamma^2\mathrm{Var}[V(s_{t+1})]+2\gamma\mathrm{Cov}(r_t,V(s_{t+1})).
\end{equation}
这直接解释了非稳态下值函数训练慢、抖动大、策略更新不稳定。

\subsection{训练-部署分布偏移的梯度失配}
策略梯度在训练分布 \(d_{\text{train}}^\pi\) 上优化：
\begin{equation}
\nabla_\theta J_{\text{train}}(\pi_\theta)=
\mathbb{E}_{(s,a)\sim d_{\text{train}}^\pi}
\left[\nabla_\theta\log\pi_\theta(a|s)\hat A_{\text{train}}(s,a)\right].
\end{equation}
若部署分布变为 \(d_{\text{test}}^\pi\)，则训练方向不再对应测试目标。常见上界为
\begin{equation}
|J_{\text{test}}(\pi)-J_{\text{train}}(\pi)|
\le
\frac{2R_{\max}}{(1-\gamma)^2}\mathrm{TV}(d_{\text{test}}^\pi,d_{\text{train}}^\pi).
\end{equation}
因此 OOD 场景中“训练分高、测试掉崖”在理论上是可预期现象。

\subsection{动作塌缩的梯度机制}
对二值动作策略 \(\pi_\theta(a=1|s)=\sigma(u_\theta(s))\)，有近似梯度形式
\begin{equation}
\frac{\partial J}{\partial u}
\propto
\mathbb{E}\left[\hat A(s,1)(1-\pi)-\hat A(s,0)\pi\right].
\end{equation}
若大量状态上 \(\hat A(s,1)\) 长期偏低，则 \(u\) 会被持续压低，导致 \(\pi(a=1|s)\to 0\)，进入自强化塌缩环。NOVA\_EDRL 的 minority 激励正是针对该机制进行反向补偿。

\subsection{阈值失配与“高 AUC 低收益”现象}
二分类决策通常写为
\begin{equation}
\hat y=\mathbf{1}[\hat p\ge \tau].
\end{equation}
在严重类别不平衡时，即使排序能力（AUC）不差，固定阈值也会导致召回崩塌，最终行为收益提升有限。该现象在 OOD 二值决策场景尤为常见。

\section{系统机理与数据协议（结合既有研究稿）}
\subsection{ALNS-RL 异步握手与共享表语义}
系统并非标准同步 Gym，而是通过共享表 \texttt{state\_reward\_pairs} 异步交互。关键哨兵值：
\begin{equation}
\texttt{action}=-10^7,\qquad \texttt{reward}=-10^7
\end{equation}
分别表示“动作未写入”和“奖励未回传”。典型流程：
\begin{enumerate}
\item ALNS 写入事件行并置空动作/奖励；
\item RL 轮询读取并写入动作；
\item ALNS 执行后回填奖励；
\item RL 在下一步读取奖励并更新。
\end{enumerate}
该异步协议导致 credit assignment 具有消息延迟特征，是本任务区别于常规 benchmark 的关键工程属性。

\subsection{阶段依赖动作语义}
动作集合虽为 \(\{0,1\}\)，但语义由阶段位 \(b_t\) 决定：
\begin{equation}
\mathrm{meaning}(a_t|b_t)=
\begin{cases}
\text{等待/保持},&a_t=0,b_t=0,\\
\text{重规划/移除},&a_t=1,b_t=0,\\
\text{接受插入},&a_t=0,b_t=1,\\
\text{拒绝插入},&a_t=1,b_t=1.
\end{cases}
\end{equation}
若不显式编码阶段上下文，会把同一动作编号的异义样本混叠，导致优势估计方向被污染。

\subsection{奖励机制的匹配本质}
默认奖励本质是“动作-可行性匹配”：
\begin{equation}
r_t=\mathbf{1}\left[(a_t=1\land \neg\mathcal{F}_t)\ \lor\ (a_t=0\land \mathcal{F}_t)\right],
\end{equation}
其中 \(\mathcal{F}_t\) 表示“当前无需重规划仍可行”。该近二值、稀疏信号会放大不平衡动作学习难度。

\subsection{分布协议：800/200 切分与模式定义}
项目中通常采用 1000 文件协议：
\begin{equation}
N=1000,\quad N_{\text{train}}=800,\quad N_{\text{impl}}=200.
\end{equation}
四类典型 pattern：
\begin{align}
\texttt{ab}:&\ A_{0:799},\ B_{800:999},\\
\texttt{aba}:&\ A_{0:399},\ B_{400:799},\ A_{800:999},\\
\texttt{abba}:&\ A_{0:399},\ B_{400:799},\ B_{800:899},\ A_{900:999},\\
\texttt{abc}:&\ A_{0:399},\ B_{400:799},\ C_{800:999}.
\end{align}
这一定义确保分布切换是“受控变量”，可用于严格比较 OOD、遗忘与泛化。

\subsection{从分布参数到策略输入的传播链}
分布影响路径可表示为：
\begin{equation}
(\text{pattern},\mu_A,\mu_B,\mu_C)
\rightarrow
\text{duration samples}
\rightarrow
\text{severity/context}
\rightarrow
o_t
\rightarrow
\pi_\theta(a_t|o_t).
\end{equation}
当上游漂移幅度大而观测压缩强时，策略输入层会出现“同观测异后果”现象，导致 OOD 误判率上升。

\section{阶段目标函数与数学定义}
\subsection{Phase1：内层预训练}
Phase1 的目标是得到稳定初始参数 \(\theta_{\text{phase1}}\)：
\begin{equation}
\theta_{\text{phase1}}=\arg\max_\theta J(\theta;E_{\text{base}}).
\end{equation}

\subsection{Phase2：固定内层，训练外层}
\subsubsection{冻结机制}
在第 \(t\) 轮，内层参数重置为同一 \(\theta_{\text{phase1}}\)，不做跨轮继承：
\begin{equation}
\theta_t \equiv \theta_{\text{phase1}},\quad \forall t\in \text{Phase2}.
\end{equation}
因此 Phase2 学到的是“如何给既有内层出有效难题”，而不是“如何继续训练内层”。

\subsubsection{Phase2 目标函数（当前实现）}
定义冻结内层回报 \(J_{\text{frozen}}(a_g)\in[0,1]\)，少数动作增益 \(G_{\text{minority}}(a_g)\)，可学习性护栏阈值 \(J_{low}\)，可行性惩罚 \(D(a_g)\ge 0\)。外层优化目标：
\begin{equation}
\mathcal{O}_{\text{P2}}(a_g)=
w_a\!\left(1-J_{\text{frozen}}(a_g)\right)
+w_m G_{\text{minority}}(a_g)
-w_{too}\!\left[\max\!\big(0,J_{low}-J_{\text{frozen}}(a_g)\big)\right]^2
-\lambda_D D(a_g).
\end{equation}

\paragraph{少数动作增益定义}
设 Phase1 内层动作占比为 \((\rho_0^{(1)},\rho_1^{(1)})\)，少数动作索引为
\begin{equation}
m = \arg\min_{k\in\{0,1\}}\rho_k^{(1)}.
\end{equation}
在给定 \(a_g\) 的冻结训练/评估中，动作占比为 \(\hat\rho_m(a_g)\)。定义
\begin{equation}
G_{\text{minority}}(a_g)=\max\left(0,\hat\rho_m(a_g)-\rho_m^{(1)}\right).
\end{equation}
直观上，若 Phase1 少数动作是 \(1\)，那么能让动作 \(1\) 占比提升的分布将得到额外奖励。

\paragraph{默认参数（代码）}
\begin{equation}
w_a=0.55,\quad w_m=0.35,\quad w_{too}=0.25,\quad J_{low}=0.15.
\end{equation}

\subsubsection{Phase2 稳定性解释}
对 \(J_{\text{frozen}}\) 求偏导（忽略不可导点的次梯度）：
\begin{equation}
\frac{\partial \mathcal{O}_{\text{P2}}}{\partial J_{\text{frozen}}}
=
-w_a
+2w_{too}\big(J_{low}-J_{\text{frozen}}\big)\mathbf{1}[J_{\text{frozen}}<J_{low}]
+w_m\frac{\partial G_{\text{minority}}}{\partial J_{\text{frozen}}}.
\end{equation}
当 \(J_{\text{frozen}}\ll J_{low}\) 时，二次护栏项增大，抑制“无限加难”；当 \(J_{\text{frozen}}>J_{low}\) 时，对抗项主导，持续寻找更具区分性的困难样本。

\subsection{Phase3：联合课程训练（保持不变）}
Phase3 恢复内层跨轮继承并保持既有 EDRL 目标：
\begin{equation}
\mathcal{O}_{\text{P3}}=
J
+\lambda_{\Delta J}\Delta J
+B(\rho_{\min})
-\lambda_D D
-\lambda_f \max(0,\rho_{floor}-\rho_{\min})
-\lambda_{fh}\max(0,\rho_{floor}-\rho_{\min})^2,
\end{equation}
其中
\begin{equation}
B(\rho_{\min})=-\eta\max(0,\rho_{target}-\rho_{\min})^{p_c}.
\end{equation}
符号说明：
\begin{itemize}
\item \(J\)：当前轮内层均值回报；
\item \(\Delta J\)：跨轮增益（相对上一轮）；
\item \(\rho_{\min}=\min(\rho_0,\rho_1)\)：动作最小占比；
\item \(D\)：路径/数据不一致惩罚；
\item \(B(\cdot)\)：偏科惩罚（低于目标占比时触发）。
\end{itemize}

\subsection{Phase4：复用 Master implement}
Phase4 不单独设计测试器，直接复用已有 implement：
\begin{equation}
\texttt{stage-mode=eval\_only},
\quad
\texttt{init-model-path}=\theta_{\text{phase3,last}},
\quad
\texttt{external-data-root}=E_{\text{phase1}}.
\end{equation}
确保与历史算法的测试口径完全一致。

\section{相位切换与收敛规则}
\subsection{Phase2 到 Phase3}
记窗口长度为 \(K\)（\texttt{converge-patience}），在最近 \(K\) 个有效轮上满足：
\begin{align}
\max_{i\in\mathcal{W}}|\Delta J_i| &\le \epsilon_{\Delta J},\\
\max_{i\in\mathcal{W}}\mathcal{O}_i-\min_{i\in\mathcal{W}}\mathcal{O}_i &\le \epsilon_O,\\
\min_{i\in\mathcal{W}}\rho_{\min,i} &\ge \rho_{\text{floor,conv}}.
\end{align}
且已达到最小迭代数 \(T_{\min}^{(2)}\) 后，切换到 Phase3；若达到 \(T_{\max}^{(2)}\) 则强制切换。

\subsection{Phase3 到 Phase4}
采用双阈值与边界控制：
\begin{itemize}
\item 连续满足稳定条件 3 次；
\item 有效迭代数满足 \(T_{\min}^{(3)}\le t \le T_{\max}^{(3)}\)；
\item 若 \(t=T_{\max}^{(3)}\) 仍未满足稳定窗，则强制切换进入 Phase4。
\end{itemize}
当前默认：\(T_{\min}^{(2)}=10, T_{\max}^{(2)}=80, T_{\min}^{(3)}=30, T_{\max}^{(3)}=200\)。

\section{外层状态、日志分解与可解释性}
\subsection{外层状态摘要}
对 \texttt{rarl\_dqn} 等策略，外层状态可由历史统计构成：
\begin{equation}
s_t^g=[J_{t-1},\Delta J_{t-1},\rho_{0,t-1},\rho_{1,t-1},\rho_{\min,t-1},\mathcal{O}_{t-1},H_{t-1},\bar J,\bar \rho_1,\bar H,\mathrm{phase},\mathrm{progress}],
\end{equation}
其中 \(H\) 为动作熵，\(\bar{\cdot}\) 为窗口均值。

\subsection{目标分解落盘}
系统会在 \texttt{outer\_actions.csv} 和 \texttt{outer\_train\_round.csv} 写入：
\begin{itemize}
\item 总目标值 \texttt{objective\_score}；
\item Phase2 分解项（如 \texttt{phase2\_term\_challenge}, \texttt{phase2\_term\_minority}, \texttt{phase2\_term\_too\_hard}, \texttt{phase2\_term\_feasibility}）；
\item 动作占比、路径读通率、样本预算使用率、停止原因等。
\end{itemize}
该机制可直接回答“外层策略在优化哪一项收益”。

\section{多 SOTA 基线矩阵与公式映射}
\subsection{RARL（ICML 2017）}
\paragraph{经典目标}
\begin{equation}
\max_{\theta}\min_{\phi}\ \mathbb{E}_{\tau\sim (\pi_\theta,\nu_\phi)}\!\left[\sum_t r_t\right].
\end{equation}
\paragraph{本系统映射}
外层作为对抗者，使用 \texttt{rarl\_dqn} 或 bandit 近似，目标口径为
\begin{equation}
\mathcal{O}_{\text{RARL}}=-J-\lambda_DD,
\end{equation}
严格零和模式下外层即时奖励 \(r_t^g=-J_t\)。

\subsection{PLR/UED（ICML 2021）}
\paragraph{核心思想}
优先重放“学习潜力高”的 level。
\paragraph{本系统映射}
定义 level 分数
\begin{equation}
s_t=w_{\Delta J}|\Delta J_t|+w_JJ_t-\lambda_DD_t.
\end{equation}
维护优先级 EMA：
\begin{equation}
\bar s_i\leftarrow (1-\alpha)\bar s_i+\alpha s_t.
\end{equation}
重放采样概率：
\begin{equation}
P(i)=\frac{\max(\varepsilon,\bar s_i-\min_j \bar s_j+\varepsilon)}{\sum_k \max(\varepsilon,\bar s_k-\min_j \bar s_j+\varepsilon)}.
\end{equation}
并采用 \texttt{new/replay} 双轨采样。

\subsection{CQL（NeurIPS 2020）}
\paragraph{经典目标}
\begin{equation}
\mathcal{L}_{\text{CQL}}
=
\mathcal{L}_{\text{TD}}
+\alpha_{\text{cql}}
\left(
\mathbb{E}_{s\sim\mathcal{D}}\left[\log\sum_a e^{Q(s,a)}\right]
-\mathbb{E}_{(s,a)\sim\mathcal{D}}[Q(s,a)]
\right).
\end{equation}
\paragraph{本系统映射}
在离散动作内层实现 \texttt{CQL\_DQN}，并输出 OOD Q 漂移相关监控量（如 \texttt{cql\_ood\_q\_gap}）。

\subsection{CaDM（ICML 2020）}
\paragraph{经典思想}
学习上下文变量 \(z\) 与动力学模型 \(p(s'|s,a,z)\) 提升跨环境泛化。
\paragraph{本系统适配}
考虑 ALNS 状态强压缩与不可微特性，采用工程化近似：在 PPO 主干增加辅助前向预测头，
\begin{equation}
\mathcal{L}_{\text{total}}
=
\mathcal{L}_{\text{PPO}}
+\beta\left(
\lambda_{\text{next}}\cdot \mathrm{MSE}(\hat y_{\text{nextsev}},y_{\text{nextsev}})
+\lambda_{\text{rew}}\cdot \mathrm{MSE}(\hat r,r)
\right).
\end{equation}
该适配保留“上下文驱动预测”思想并兼顾可训练性。

\subsection{SOTA 资源本地化路径}
论文与代码已下载至：
\begin{itemize}
\item \texttt{ALNS\_Research\_Documentation/SOTA\_Baselines/RARL}
\item \texttt{ALNS\_Research\_Documentation/SOTA\_Baselines/PLR\_UED}
\item \texttt{ALNS\_Research\_Documentation/SOTA\_Baselines/CQL}
\item \texttt{ALNS\_Research\_Documentation/SOTA\_Baselines/CADM}
\end{itemize}

\section{评测协议与关键指标}
\subsection{分布矩阵}
主评测覆盖 OOD、Generalization、Forgetting、Random Mix 四组。建议主文采用 16 个黄金分布作为报告主矩阵。

\subsection{Normalized Policy Score（NPS）}
对测试分布 \(\mathcal{D}\)，定义
\begin{equation}
\mathrm{NPS}(\pi_{RL},\mathcal{D})
=
\frac{J(\pi_{RL},\mathcal{D})-J(\pi_{rand},\mathcal{D})}
{\max\!\left(J(\pi_{a=0},\mathcal{D}),J(\pi_{a=1},\mathcal{D})\right)-J(\pi_{rand},\mathcal{D})}.
\end{equation}
其中 \(J(\cdot,\mathcal{D})\) 使用 implement 阶段样本均值估计。

NPS 的解释：
\begin{itemize}
\item \(\mathrm{NPS}\le 0\)：不优于随机策略；
\item \(0<\mathrm{NPS}<1\)：优于随机但未超越最佳静态规则；
\item \(\mathrm{NPS}=1\)：达到最佳静态规则水平；
\item \(\mathrm{NPS}>1\)：实现上下文自适应超越。
\end{itemize}

\section{算法流程（可复现实验口径）}
\begin{algorithm}[H]
\caption{NOVA\_EDRL 四阶段训练与评测}
\begin{algorithmic}[1]
\State 输入：基础分布配置、\(\texttt{run\_id}\)、\(\texttt{base\_seed}\)、预算与收敛阈值
\State \textbf{Phase1}：训练内层，得到 \(\theta_{\text{phase1}}\)
\State 初始化外层策略参数 \(\psi\)
\For{\(t=1,2,\dots\)}
  \State 依据 \(s_t^g\) 选择外层动作 \(a_{g,t}\)
  \State 由 \(h(\texttt{base\_seed},t)\) 生成确定性 seed，并生成小批量事件文件
  \If{当前为 Phase2}
    \State 内层从 \(\theta_{\text{phase1}}\) 启动并运行预算步
    \State 计算 \(\mathcal{O}_{\text{P2}}\)，更新 \(\psi\)
    \If{满足 Phase2 收敛或达到上限} \State 切换至 Phase3 \EndIf
  \Else
    \State 内层跨轮继承更新，计算 \(\mathcal{O}_{\text{P3}}\)，更新 \(\psi\)
    \If{满足 Phase3 收敛或达到上限} \State \textbf{break} \EndIf
  \EndIf
\EndFor
\State \textbf{Phase4}：调用 Master implement（\texttt{eval\_only}）输出最终评测
\end{algorithmic}
\end{algorithm}

\section{工程实现映射}
\begin{table}[H]
\centering
\caption{NOVA\_EDRL 与 SOTA 矩阵的关键实现模块}
\begin{tabular}{p{4.9cm}p{9.5cm}}
\toprule
模块 & 作用 \\
\midrule
\texttt{codes/Dynamic\_master34959.py} & 算法入口；支持 NOVA\_EDRL、RARL、PLR\_UED、CQL\_DQN、CADM \\
\texttt{codes/outer\_rl/run\_edrl\_pipeline.py} & 四阶段总编排；Phase4 复用 Master implement \\
\texttt{codes/outer\_rl/run\_edrl\_phase2.py} & 外层动作采样、目标计算、切换判据、日志分解 \\
\texttt{codes/generation/generate\_outer\_small\_batch.py} & 外层动作到事件文件批次的映射生成 \\
\texttt{codes/core/dynamic\_RL34959.py} & 内层 RL 主循环与 ALNS 联动 \\
\texttt{codes/robust\_rl/cql\_dqn.py} & CQL\_DQN 离散保守学习 \\
\texttt{codes/robust\_rl/ppo\_new/v6\_cadm.py} & CaDM 适配版 PPO（辅助预测损失） \\
\texttt{codes/analysis/compute\_metrics.py} & 汇总指标计算（含 NPS） \\
\bottomrule
\end{tabular}
\end{table}

\section{实验结果（待填写）}
\textcolor{gray}{本节保留用于填入正式实验结果：主结果表、消融、收敛曲线、跨分布泛化、NPS 统计显著性分析。}

\section{结论（待填写）}
\textcolor{gray}{本节暂不写结论，待全部实验完成后统一撰写。}

\section*{附录 A：当前默认超参数}
\begin{itemize}
\item 外层默认策略：\texttt{policy-mode=ts}，目标模式 \texttt{objective-mode=edrl}
\item Phase2 迭代窗：\([10,80]\)；Phase3 迭代窗：\([30,200]\)
\item Phase2 权重：\(w_a=0.55,\ w_m=0.35,\ w_{too}=0.25,\ J_{low}=0.15\)
\item 路径惩罚系数：\(\lambda_D=2.0\)
\item 收敛窗口：\(\texttt{patience}=3,\ \epsilon_{\Delta J}=0.05,\ \epsilon_O=0.20\)
\end{itemize}

\section*{附录 B：SOTA 对照的本地索引}
\begin{itemize}
\item \texttt{ALNS\_Research\_Documentation/SOTA\_Baselines/RARL/papers}
\item \texttt{ALNS\_Research\_Documentation/SOTA\_Baselines/PLR\_UED/papers}
\item \texttt{ALNS\_Research\_Documentation/SOTA\_Baselines/CQL/papers}
\item \texttt{ALNS\_Research\_Documentation/SOTA\_Baselines/CADM/papers}
\end{itemize}

\end{document}
